\documentclass{article}
\usepackage[letterpaper, textwidth=5.5in, textheight=9.0in]{geometry}
\usepackage{graphicx} % Required for inserting images

% text and typesetting macros
\usepackage{setspace}
\setstretch{1.08}
\renewcommand{\paragraph}[1]{\medskip\par\noindent{\textbf{#1}}~---}
\sloppy\sloppypar\raggedbottom
\frenchspacing
\newcommand{\documentname}{\textsl{Note}}
\newcommand{\sectionname}{Section}
\newcommand{\secref}[1]{\sectionname~\ref{#1}}

% math macros
\usepackage{amsmath}
\newcommand{\e}{\mathrm{e}}
\newcommand{\given}{\,|\,}

% collaboration / writing / commenting macros
\usepackage{xcolor}
\newcommand{\hoggsay}[1]{{\color{blue} (Hogg says: #1)}}

\title{\bfseries A custom periodogram for \textsl{Gaia} astrometric residuals.}
\author{Hogg, Casey for Ramirez~Delgado and friends}
\date{\textit{draft version ---} 2026 October 20}
\begin{document}
\maketitle

\paragraph{Abstract}
We provide a replacement for the periodogram that knows about ``scan angle'' and also knows that the signals we care about are elliptically polarized.

\section{Introduction}\label{sec:intro}
ESA \textsl{Gaia}'s data is fundamentally one-dimensional, and it is entirely in the time domain:
To a very good approximation, \textsl{Gaia} scans the sky and measures the \emph{times} at which stars transit a fiducial point in each of two focal planes.
In this brief \documentname, we ask:
\begin{itemize}
    \item If \textsl{Gaia} does a good job of fitting every object with some kind of kinematic model (parallax and proper motion, say, or even a more sophisticated model), and
    \item if \textsl{Gaia} releases the timing residuals away from that model (or code to compute them), then
    \item can we find periodic signals (of the kinds relevant to astronomy) with some kind of modified periodogram?
\end{itemize}
We are writing this note, so obviously the answer is ``yes.''

\section{Assumptions and choices}\label{sec:assumptions}
\paragraph{Residual data}
For the purposes of this \documentname, we will assume that a very good kinematic model of every star has been subtracted, and we are left with good residuals.
The model could be generalized, of course.
The data we are envisioning (and this may get confusing) is that we have $N$ timing residuals $\Delta_n$ ($1\leq n\leq N$).
Each of these residuals is measured at an observation time $t_n$ and has an uncertainty variance $\sigma^2_n$.
These timing residuals $\Delta_n$ have units of time (perhaps in seconds) and they happen at times $t_n$ (perhaps also in seconds).
Importantly the times $t_n$ will be some kind of (Solar System) barycentric times.

Except for certain weird modes, Gaia scans at a constant angular scan rate $\Omega$ such that the astrometric residual $y_n$ corresponding to timing residual $\Delta_n$ is just $y_n = \Omega\,\Delta_n$.
In what follows, we will mainly work with $y_n$ not $\Delta_n$ because we care most about quantities on the sky (not quantities on the spacecraft).
% ARC: I think Gaia defines psi north towards east
Most importantly for what follows, at each observation time $t_n$, the spacecraft is scanning the sky at scan angle $\psi_n$ in the plane of the sky, defined relative to some fiducial two-dimensional on-sky coordinate system.

\paragraph{Circular orbits projected on the sky}
The (extremely restrictive) model we have in mind is that the residuals are formed by some kind of periodic circular motion of the photocenter of the star.
Projected on the sky, this becomes periodic elliptical motion.
Periodic elliptical motion can be described by
\begin{align}
    \vec{\theta}(t) &= \vec{a}\,\cos(\omega\,t-\phi_0) + \vec{b}\,\sin(\omega\,t-\phi_0) \\
    0 &= \vec{a}\cdot\vec{b} ~,
\end{align}
where $\vec{\theta}$ is a 2-vector offset in the plane of the sky (or, if you like, on a relevant tangent plane),
$\vec{a}$ and $\vec{b}$ are the semi-major and semi-minor axis vectors on the plane of the sky,
$\omega$ is the angular frequency of the circular motion,
and $\phi_0$ is a phase corresponding to maximum elongation.
Note that the sign of $\vec{a}\times\vec{b}\equiv\vec{a}\wedge\vec{b}$ (which, in 2 dimensions, is a pseudo-scalar) defines a parity for the orbit on the sky.

\paragraph{Generative model}
The fundamental assumption is that we see noisy projections of the periodic elliptical offfset motion.
That is,
\begin{align}
    y_n &= \hat{\psi}_n\cdot\vec{\theta}(t_n) + s_n \label{eq:genmodel}\\
    \hat{\psi}_n &= \hat{x}\,\cos\psi_n + \hat{y}\,\sin\psi_n ~,
\end{align}
where $\hat{\psi}_n$ is a unit 2-vector defining pointing in the scan direction,
$s_n$ is a ``residual'' or ``noise'' or ``incorrectness'' (we will call it the latter so as not to overload the word ``residual'' and not to attribute too much to the meaning of the $s_n$),
and $\hat{x},\hat{y}$ are orthonormal 2-vectors defining the on-sky coordinate system in which the scan angle $\psi_n$ is defined (and indeed in which all of our on-sky 2-vectors are defined).

All of these expressions, by the way, depend on the handedness of the \textsl{Gaia} definition of the scan angle $\psi$.
In some sense, the \textsl{Gaia} sign convention sets the handedness of the $\hat{x},\hat{y}$ coordinate system,
and the meaning of a positive value for $\vec{a}\wedge\vec{b}$.
That is, there is a correspondence between the handedness of the scan-angle definition, the handedness of the $\hat{x},\hat{y}$ coordinate system,
the handedness of the $\vec{a},\vec{b}$ direction pair, and the elliptical polarization of the residual motion we are modeling.
For our own scientific purposes, we don't care about this overall sign, but for some purposes it might be important to understand or control this.

We can rotate the expression \eqref{eq:genmodel} into a frame that is aligned with the semi-axes $\vec{a},\vec{b}$ and take the dot product explicitly to get
\begin{align}
    y_n &= |\vec{a}|\,\cos(\psi_n-\psi_0)\,\cos(\omega\,t-\phi_0) \pm |\vec{b}|\,\sin(\psi_n-\psi_0)\,\sin(\omega\,t-\phi_0) ~,
\end{align}
where $\psi_0$ is the scan angle that aligns the scan direction with the semi-major axis vector $\vec{a}$,
and the sign of the $|\vec{b}|$ term is related to the parity (the sign of $\vec{a}\wedge\vec{b}$). \hoggsay{This whole section needs a good diagram on the plane of the sky!}

\paragraph{Noise and robustness}
We will assume that the incorrectnesses $s_n$ are are drawn from a noise distribution $p(s\given\sigma^2)$ controlled by a variance parameter $\sigma^2$, such that each $s_n$ is drawn from $p(s_n\given\sigma^2_n)$.
In the simplest case this distribution would be the zero-mean Gaussian.
In more robust inferences we might make this something with much fatter tails than that.

\section{Method}\label{sec:method}
If you carry forward the mathematics, and apply either useful trig identities, or else complex exponentials, you get to
\begin{align}
    y_n = Z_+\,\e^{i\,[\omega\,t_n + \psi_n]} + Z_-\,\e^{i\,[\omega\,t_n - \psi_n]} + s_n ~, \label{eq:zs}
\end{align}
where really we take the real part of this expression,
and $Z_+, Z_-$ are complex coefficients.
You might ask why we have two amplitudes $Z_+, Z_-$ and not just one; it's because these amplitudes parameterize not just the amplitude and phase of the oscillating signal; they also parameterize a polarization tensor on the sky.
That is, we want to know the complex coefficients; the magnitudes of the coefficients will relate to the semi-amplitudes $|\vec{a}|,|\vec{b}|$; the ratio of imaginary to real parts of the coefficients will relate
to the offsets $\phi_0,\psi_0$.
In detail, the complex coefficients are related to the vectors and angles and phases as follows:
\begin{align}
    |\vec{a}| &= |Z_+|+|Z_-| \\
    |\vec{b}| &= ||Z_+|-|Z_-|| \\
    \phi_0 &= -\tfrac{1}{2}\,\bigl[\arg Z_+ + \arg Z_-\bigr] \\
    \psi_0 &= +\tfrac{1}{2}\,\bigl[\arg Z_- - \arg Z_+\bigr] ~,
\end{align}
where $\arg Z \equiv \operatorname{atan2}(\operatorname{Im} Z, \operatorname{Re} Z)$ is the complex phase.

The simplest (and slowest) version of the method we are proposing here is to use \eqref{eq:zs} to build a design matrix $X$ and then do
\begin{align}
    \hat{\beta} &= [X^\top\,C^{-1}\,X]^{-1}\,X^\top\,C^{-1}\,Y ~,
\end{align}
where
\begin{align}
    \beta &\equiv \begin{bmatrix} c_+ & s_+ & c_- & s_- \end{bmatrix}^\top ~,
    \quad\mbox{with}\quad Z_\pm = c_\pm - i\,s_\pm ~,
\end{align}
and the design matrix is shape $N \times 4$ with entries
\begin{align}
    X_n &= \begin{bmatrix}
        \cos(\omega\,t_n + \psi_n) &
        \sin(\omega\,t_n + \psi_n) &
        \cos(\omega\,t_n - \psi_n) &
        \sin(\omega\,t_n - \psi_n)
    \end{bmatrix} ~,
\end{align}
This is the pure-real implementation of this (4 real parameters);
in the most naive complex implementation (2 complex parameters)
both the design matrix $X$ and the output $[Z_+,Z_-]$ are complex in the most naive implementation of this.

\hoggsay{A more robust method is to do what now?}

\hoggsay{A much faster method is to do what now?}

\paragraph{Implementation notes}
\hoggsay{What did we actually do?}

\section{Experiments and results}\label{sec:experiments}

\section{Discussion}\label{sec:discussion}

\paragraph{Acknowledgments}
It is a pleasure to thank Adrian Price-Whelan (Flatiron) and \hoggsay{other people}
for help with all this.

\end{document}
